%%%%%
%
%% Dans ce fichier, les �ventuels accents sont cod�s
%% en ISO-8859-15 (aka latin9)
%% La ligne suivante donne trois exemples 
%% e accent aigu : �, e accent grave : �, e-dans-l'o : �
\section{Hauteur de la base des nuages moyenne des 3 couches de nuages quand le nombre de couches reste constant}

\itemstyle{Op\'erateur~:} \operfmt{HbnMoyInvEch\_tps\_bnds}
\index{HbnMoyInvEch\_tps\_bnds}

%%%%%OBJET%%%%
\itemstyle{Objet~:}
Cette fonction calcule, pour chaque instant $t$ d'une
base de temps $time_{out}$, 
la moyenne de chaque couche de hauteur de base de nuage $X_{i_{out}}$ entre les bornes donn�es par le $time_{bnds}$ associ� � $time_{out}$.
Le calcul n'est r�alis� que pour des valeurs < $val_{inv}$ et si le nombre de couches < $val_{inv}$ reste constant
Aucune supposition n'est faite sur la r�gularit� de
l'�chantillonnage des 2 variables. Il est cependant n�cessaire
que $time_{in}$ (le temps des $X_{i_{in}}$) soit "� peu pr�s" tri� par ordre
croissant.


%%%%%ARGUMENTS%%%%
\bigskip
%% Dans ce fichier, les �ventuels accents sont cod�s
%% en ISO-8859-15 (aka latin9)
%% La ligne suivante donne trois exemples 
%% e accent aigu : �, e accent grave : �, e-dans-l'o : �
%%%%%ENTREES%%%%
\begin{tabarg}
\multicolumn{3}{l}{\titretab {Entr\'ees}}\\
\hline
$X1_{in}$ & Vecteur, \texttt{float} & Hauteur de la base des nuages de la premi�re couche\\
$X2_{in}$ & Vecteur, \texttt{float} & Hauteur de la base des nuages de la deuxi�me couche\\
$X3_{in}$ & Vecteur, \texttt{float} & Hauteur de la base des nuages de la troisi�me couche\\
$toler_{mnq}$ & Coeff., \texttt{Int\_32} & Seuil de calcul de la moyenne (nb. �chant). A chaque instant de la base de temps d'arriv�e, 
la moyenne est calcul�e si le nombre d'�chantillons dans l'intervalle de temps est strictement sup�rieur � $toler_{mnq}$. Sinon, 
le r�sultat est "valeur manquante". Si $toler_{mnq}$ = 0, le calcul est fait d�s qu'il y a au moins une valeur valide dans l'�chantillon. \\
$time_{out}$ & Vecteur, \texttt{double} & Base de temps d'arriv�e (secondes)\\
$val_{inv}$ & Coeff., \texttt{float} & valeur invalide pour laquelle la Moyenne n'est pas calcul�e \\
$val_{remplac}$ & Coeff., \texttt{float} & valeur de remplacement quand la valeur atteint val\_{inv} avec une tol�rance \\
$tolerance_{inv}$ & Coeff., \texttt{float} & proportion de valeurs n�cessaires pour que le calcul de la moyenne soit valide (entre 0 et 1). Si le nb de valeurs utilis�es pour calculer la moyenne divis� par le nombre de valeurs de la p�riode est < tol�rance\_{inv}, la valeur de la moyenne est mise �  val\_{remplac} \\
\end{tabarg}

\medskip
%%%%%SORTIES%%%%
\begin{tabarg}
\multicolumn{3}{l}{\titretab {Sorties}}\\
\hline
$X1_{out}$ & Vecteur, \texttt{float} & $X1_{in}$ moyenn� et �chantillonn� sur
la base de $time_{out}$ (m�me unit� que $X1_{in})$ si le nb de couches < $val_{inv}$ reste constant\\
$X2_{out}$ & Vecteur, \texttt{float} & $X2_{in}$ moyenn� et �chantillonn� sur
la base de $time_{out}$ (m�me unit� que $X2_{in})$ si le nb de couches < $val_{inv}$ reste constant\\
$X3_{out}$ & Vecteur, \texttt{float} & $X3_{in}$ moyenn� et �chantillonn� sur
la base de $time_{out}$ (m�me unit� que $X3_{in})$ si le nb de couches < $val_{inv}$ reste constant\\
\end{tabarg}



%%%%%FORMULE%%%%
\itemstyle{Formule}

%% Formule de l'operateur : HbnMoyInvEch_tps_bnds
%% Dans ce fichier, les �ventuels accents sont cod�s
%% en ISO-8859-15 (aka latin9)
%% La ligne suivante donne trois exemples 
%% e accent aigu : �, e accent grave : �, e-dans-l'o : �

Pour chaque couche Xn:

Pour  i=$0$, Nbre  $Time_{bnds}$ :

\[
Xn_{out}[i] =
\begin{cases}
\begin{tabular}{ll}
$\frac{\displaystyle{\sum_{tb[i][0]}^{tb[i][1]} Xn_{in}[j]} }{Nn_i}$
&
\texttt{ si  nb couche inf\'erieure \`a $val_{inv}$ est constant}
\\
\texttt{$val_{replac}$} & \texttt{ sinon}
\end{tabular} 
\\
\end{cases}
\]

Avec $Xn_{in}[j]$ inf\'erieur \`a $val_{inv}$ et  $Nn_i$ le nombre de donn\'ees dans l'intervalle tb  pour lesquels $Xn_{in}[j]$ est inf\'erieur \`a $val_{inv}$ .


%%%%COMMENTAIRE%%%%
\itemstyle{Commentaire~:}
\textit{tb} est un tableau d'intervalles de temps d�termin�s � partir du vecteur temps $time_{out}$.

%%%%%AUTEUR%%%%
\itemstyle{Auteur~:}
CNRM/GMEI/TRAMM

%%%%%REFERENCE%%%%
\itemstyle{R\'ef\'erences~:}
N�ant.
